% Generated draft: requires mathematical and source review.
\documentclass[11pt]{article}
\usepackage[margin=25mm]{geometry}
\usepackage{fontspec}
\setmainfont{Latin Modern Roman}
\usepackage{amsmath,amssymb,mathtools,unicode-math}
\setmathfont{Latin Modern Math}
\usepackage{enumitem,needspace,xurl,hyperref}
\setlength{\parindent}{0pt}
\setlength{\parskip}{5pt}
\setlength{\emergencystretch}{3em}
\newcommand{\N}{\mathbb{N}}
\newcommand{\Z}{\mathbb{Z}}
\newcommand{\Q}{\mathbb{Q}}
\newcommand{\R}{\mathbb{R}}
\newcommand{\C}{\mathbb{C}}
\newcommand{\F}{\mathbb{F}}
\newcommand{\A}{\mathbb{A}}
\newcommand{\PP}{\mathbb{P}}
\newcommand{\E}{\mathbb{E}}
\newcommand{\eps}{\varepsilon}
\newcommand{\dd}{\,\mathrm{d}}
\begin{document}
\section*{UnsolvedMath ID: 30003685}
\textbf{Source problem code:} catalogue entry
\subsection{Definitions and mathematical statement}
The supplied record poses the historical classification question:
\begin{quote}
Were early modern optical diagrams mathematical instruments?
\end{quote}
The terms ``early modern optical diagram'' and ``mathematical instrument'' are not defined in the supplied record. In particular, no corpus of diagrams, chronological boundaries for ``early modern'', or criterion for classification as a mathematical instrument is specified. The supplied assessment describes the definition as partly arbitrary and supplies no settled criterion resolving the question. The question also does not specify whether it concerns all such diagrams, some examples, or their classification in particular historical contexts.

The supplied statement-verification note identifies this as a historiographical classification question rather than a mathematical problem, marks the statement as \texttt{unrecoverable}, and states that no separate clean statement is promoted; the original record remains publicly visible for provenance. Consequently, the question is preserved here without inventing mathematical definitions, hypotheses, or quantifiers. A precise mathematical proposition cannot be recovered from the supplied evidence.
\subsection{Short English statement}
The question asks whether early modern optical diagrams should be classified as mathematical instruments. The supplied record leaves the relevant terms, historical scope, and classification criteria unresolved, so it supports a historical question rather than a precise mathematical statement. Its supplied ``open'' label is a curation label, not an independently verified mathematical status.
\subsection{Sources}
\begin{itemize}
\item Local UnsolvedMath registry, permanent ID \texttt{30003685}, source problem code \texttt{catalogue entry}, title ``Historical Classification of Optical Diagrams as Instruments''. The supplied record belongs to the Oberwolfach Reports Open Problems collection, records proposed year 2017, and names no proposer.
\item Supplied source citation: \emph{Mathematical Instruments between Material Artifacts and Ideal Machines: Their Scientific and Social Role before 1950} (2018), \emph{Oberwolfach Reports}, DOI 10.4171/owr/2017/58, pp.~3471--3560. \url{https://doi.org/10.4171/owr/2017/58}. The supplied background describes these reports as documenting mathematical workshops and their open-problem discussions; it does not supply a passage establishing precise classification criteria for this question.
\item Supplied literature source: ``Oberwolfach report'', \url{https://doi.org/10.4171/owr/2017/58}. Literature checked at: \texttt{2026-08-21}. This is supplied dated curation, not an independently verified assessment. It labels the question ``open'' and states that the definition is partly arbitrary and that no settled criterion resolving the definitional question is supplied. The background requests independent verification of time-sensitive status and cited claims.
\end{itemize}
\end{document}
